\section{Physics CFD: Version 1 Junction Stokes Library}

\subsection{Purpose and Scientific Scope}

The CFD module provides a spatially resolved, local flow model for the
validated upper junction of the asymmetric loop device. It converts a prescribed
branch-flow split into a continuous velocity and pressure field on the junction
patch. In Version 1, the branch split is an input, not a prediction:

\[
  f_L = \frac{Q_L}{Q_{\mathrm{in}}},
  \qquad
  f_R = 1 - f_L.
\]

The CFD solver computes \(\mathbf{u}(\mathbf{x})\) and \(p(\mathbf{x})\) for a
given pair \((f_L, f_R)\). The future hydraulics module is expected to provide
the split used by the CFD module. The transformer should consume CFD-derived
fields or features only after the physics pipeline has selected the split.

Version 1 intentionally does not implement:

\begin{itemize}
  \item prediction of branch split from pressure or droplet state,
  \item interpolation between CFD split cases,
  \item transient CFD,
  \item droplet-resolved multiphase CFD,
  \item transformer integration.
\end{itemize}

\subsection{Files}

\begin{itemize}
  \item \texttt{src/physics/cfd/domain.py}: CFD-domain construction from the
        validated centerline and channel width.
  \item \texttt{src/physics/cfd/mesh.py}: production triangular mesh generation,
        mesh-quality reporting, topology diagnostics, and mesh figures.
  \item \texttt{src/physics/cfd/solver.py}: steady Stokes P2/P1 finite-element
        solve, boundary conditions, metadata, and field/streamline outputs.
  \item \texttt{src/physics/cfd/inlet\_profile\_diagnostic.py}: raw-FEM inlet
        Poiseuille validation.
  \item \texttt{src/physics/cfd/mesh\_convergence.py}: focused default-vs-fine
        mesh comparison.
  \item \texttt{src/physics/cfd/pilot\_splits.py}: three pilot prescribed-split
        cases used before the full library.
  \item \texttt{src/physics/cfd/velocity\_library.py}: full Version 1 prescribed
        split library generation and indexing.
  \item \texttt{scripts/physics/build\_junction\_mesh.py}: production mesh
        generation command.
  \item \texttt{scripts/physics/solve\_junction\_cfd.py}: single prescribed
        split solve command.
  \item \texttt{scripts/physics/compare\_default\_fine\_mesh.py}: focused mesh
        convergence command.
  \item \texttt{scripts/physics/run\_pilot\_splits.py}: pilot split runner.
  \item \texttt{scripts/physics/generate\_velocity\_library.py}: full library
        batch runner.
  \item \texttt{configs/physics/junction\_cfd.yml}: production CFD
        configuration.
\end{itemize}

\subsection{Physical Geometry and Units}

The CFD domain is the reduced upper-junction window of the asymmetric-loop
device. The geometry is constructed from the validated centerline and the known
channel width, not from pixelated segmentation masks. The physical channel
width is

\[
  W = 100~\mu\mathrm{m},
\]

and the calibration is

\[
  1~\mathrm{px} = 4~\mu\mathrm{m}.
\]

All geometry coordinates are stored in micrometers. The finite-element solve is
assembled in SI-compatible units by converting coordinates to meters:

\[
  \mathbf{x}_{m} = 10^{-6}\mathbf{x}_{\mu m}.
\]

The reported field units are:

\begin{center}
\begin{tabular}{@{}ll@{}}
\toprule
Quantity & Unit \\
\midrule
Geometry coordinates & \(\mu\mathrm{m}\) \\
Solver coordinates & \(\mathrm{m}\) \\
Velocity & \(\mathrm{m/s}\) \\
Pressure & \(\mathrm{Pa}\) \\
Dynamic viscosity & \(\mathrm{Pa\,s}\) \\
Two-dimensional flux & \(\mathrm{m^2/s}\) \\
\bottomrule
\end{tabular}
\end{center}

\subsection{Governing Equations}

Version 1 solves the steady incompressible Stokes equations for a Newtonian
single-phase flow:

\[
  -\nabla p + \mu \nabla^2 \mathbf{u} = \mathbf{0},
  \qquad
  \nabla \cdot \mathbf{u} = 0.
\]

Here \(\mathbf{u} = (u_x,u_y)\) is the in-plane velocity, \(p\) is pressure, and
\(\mu\) is dynamic viscosity. In weak form, the viscous and incompressibility
terms are assembled as

\[
  a(\mathbf{u},\mathbf{v})
  =
  \int_\Omega \mu \nabla \mathbf{u} : \nabla \mathbf{v}\,d\Omega,
\]

\[
  b(\mathbf{u},q)
  =
  -\int_\Omega q\,\nabla\cdot\mathbf{u}\,d\Omega.
\]

The resulting mixed saddle-point system has the block form

\[
  \begin{bmatrix}
  A & B^T \\
  B & 0
  \end{bmatrix}
  \begin{bmatrix}
  \mathbf{U} \\
  \mathbf{P}
  \end{bmatrix}
  =
  \begin{bmatrix}
  \mathbf{f} \\
  \mathbf{0}
  \end{bmatrix},
\]

with essential velocity boundary conditions imposed by condensation.

\subsection{Boundary Conditions}

The inlet velocity is prescribed as a planar Poiseuille profile:

\[
  u_{\parallel}(s)
  =
  \frac{3}{2}U_{\mathrm{mean}}
  \left[
  1 - \left(\frac{2s}{W}\right)^2
  \right],
  \qquad
  -\frac{W}{2}\le s \le \frac{W}{2}.
\]

The mean inlet velocity is obtained from the experimental total flow rate by
converting the three-dimensional volumetric flow rate to a two-dimensional flux
per unit depth:

\[
  q_{\mathrm{2D}} = \frac{Q_{\mathrm{3D}}}{H},
  \qquad
  U_{\mathrm{mean}} = \frac{q_{\mathrm{2D}}}{W}.
\]

For the configured case,

\[
  Q_{\mathrm{total}}=1960~\mu\mathrm{L/hr},
  \qquad
  H=100~\mu\mathrm{m},
  \qquad
  q_{\mathrm{2D}} = 5.444444\times10^{-6}~\mathrm{m^2/s}.
\]

The outlet velocities are also prescribed as parabolic profiles. For a requested
left fraction \(f_L\),

\[
  q_L = f_L q_{\mathrm{2D}},
  \qquad
  q_R = (1-f_L)q_{\mathrm{2D}}.
\]

The wall boundary condition is no slip:

\[
  \mathbf{u}=\mathbf{0}
  \qquad
  \mathrm{on~solid~walls}.
\]

No physical pressure Dirichlet condition is applied at the inlet or outlets.
Exactly one P1 pressure degree of freedom is fixed as a gauge to remove the
constant pressure null mode. A sparse direct solve is required. Singular systems
fail loudly; no least-squares fallback is accepted.

\subsection{Finite Elements and Degree-of-Freedom Handling}

The solver uses the Taylor--Hood element pair:

\[
  \mathbf{u}\in [P_2]^2,
  \qquad
  p\in P_1.
\]

Velocity degrees of freedom are vector-valued P2 DOFs, including edge-midpoint
DOFs. Pressure degrees of freedom are scalar P1 nodal DOFs. The solution vector
stores all velocity coefficients first, followed by pressure coefficients.
Velocity boundary conditions are applied only on actual boundary-facet trace
DOFs: edge endpoints and P2 edge midpoints. This was an important Version 1
stability fix. Earlier tolerance-band boundary selection accidentally
constrained near-boundary interior velocity DOFs and produced singular mixed
systems. The current formulation reports full structural rank for the condensed
system.

\subsection{Production Mesh}

The production mesh is a triangular mesh generated from centerline-derived
boundary samples and a filtered interior point cloud. Boundary samples are
protected first; later candidate points closer than \(6~\mu\mathrm{m}\) are
rejected. This removes near-duplicate samples that previously caused sliver
triangles without changing the physical domain.

\begin{center}
\begin{tabular}{@{}lr@{}}
\toprule
Production mesh metric & Value \\
\midrule
Nodes & 425 \\
Elements & 656 \\
Minimum angle & \(20.299^\circ\) \\
Mean aspect ratio & 1.400 \\
Maximum aspect ratio & 2.883 \\
Connected fluid components & 1 \\
Domain holes & 0 \\
Invalid/inverted elements & 0 \\
Interior boundary facets & 0 \\
\bottomrule
\end{tabular}
\end{center}

\begin{figure}[H]
  \centering
  \includegraphics[width=0.9\linewidth,height=0.70\textheight,keepaspectratio]{../outputs/physics/junction_cfd/figures/junction_mesh.png}
  \caption{Production triangular CFD mesh for the reduced upper-junction domain.}
  \label{fig:cfd-production-mesh}
\end{figure}

\begin{figure}[H]
  \centering
  \includegraphics[width=0.9\linewidth,height=0.70\textheight,keepaspectratio]{../outputs/physics/junction_cfd/figures/junction_mesh_zoom.png}
  \caption{Zoomed production mesh near the bifurcation. The sharp junction geometry is preserved.}
  \label{fig:cfd-production-mesh-zoom}
\end{figure}

\begin{figure}[H]
  \centering
  \includegraphics[width=0.9\linewidth,height=0.70\textheight,keepaspectratio]{../outputs/physics/junction_cfd/figures/mesh_quality_aspect_ratio.png}
  \caption{Production mesh quality visualization colored by triangle aspect ratio.}
  \label{fig:cfd-production-mesh-quality}
\end{figure}

\subsection{Fine-Mesh Evidence}

A focused convergence check compared the production mesh against one finer mesh.
The finer mesh uses the same geometry, boundary labels, solver formulation, and
boundary conditions. It is retained as convergence evidence only and is not the
default.

\begin{center}
\begin{tabular}{@{}lrr@{}}
\toprule
Metric & Production & Fine evidence mesh \\
\midrule
Nodes & 425 & 971 \\
Elements & 656 & 1619 \\
Minimum angle & \(20.299^\circ\) & \(19.081^\circ\) \\
Mean aspect ratio & 1.400 & 1.245 \\
Maximum aspect ratio & 2.883 & 2.824 \\
\bottomrule
\end{tabular}
\end{center}

On common physical sample points in the junction region, the mean
velocity-magnitude difference was \(6.062292\times 10^{-4}~\mathrm{m/s}\), and
the mean flow-direction difference was \(0.909839^\circ\). The seeded streamline
topology remained \(15\) left, \(15\) right, and one center separatrix/other for
the 50/50 case. The production mesh was therefore retained as the default.

\subsection{Validated 50/50 Case}

The validated reference case uses

\[
  f_L = f_R = 0.5.
\]

\begin{center}
\begin{tabular}{@{}lr@{}}
\toprule
Quantity & Value \\
\midrule
Solver backend & \texttt{scikit-fem/direct} \\
Measured left/right split & 0.5000 / 0.5000 \\
Inlet flux & \(-5.444444\times10^{-6}~\mathrm{m^2/s}\) \\
Left outlet flux & \(2.722222\times10^{-6}~\mathrm{m^2/s}\) \\
Right outlet flux & \(2.722222\times10^{-6}~\mathrm{m^2/s}\) \\
Net mass residual & \(-1.730446\times10^{-17}~\mathrm{m^2/s}\) \\
Maximum velocity & \(8.166667\times10^{-2}~\mathrm{m/s}\) \\
Pressure range & \(-2.409284\) to \(26.29512~\mathrm{Pa}\) \\
Seeded streamlines & 15 left, 15 right, 1 separatrix/other \\
\bottomrule
\end{tabular}
\end{center}

\begin{figure}[H]
  \centering
  \includegraphics[width=0.9\linewidth,height=0.70\textheight,keepaspectratio]{../outputs/physics/junction_cfd/solutions/split_0p50/figures/velocity_magnitude.png}
  \caption{Velocity magnitude for the validated 50/50 prescribed-split Stokes solution.}
  \label{fig:cfd-50-velocity}
\end{figure}

\begin{figure}[H]
  \centering
  \includegraphics[width=0.9\linewidth,height=0.70\textheight,keepaspectratio]{../outputs/physics/junction_cfd/solutions/split_0p50/figures/pressure_field.png}
  \caption{Pressure field for the validated 50/50 solution. Pressure is defined up to the one-DOF gauge.}
  \label{fig:cfd-50-pressure}
\end{figure}

\begin{figure}[H]
  \centering
  \includegraphics[width=0.9\linewidth,height=0.70\textheight,keepaspectratio]{../outputs/physics/junction_cfd/solutions/split_0p50/figures/velocity_streamlines_inlet_seeded.png}
  \caption{Inlet-seeded streamlines for the 50/50 case. Streamlines are post-processing traces of the solved velocity field.}
  \label{fig:cfd-50-seeded-streamlines}
\end{figure}

\subsection{Inlet Poiseuille Validation}

The raw FEM field is sampled directly at three inlet cross-sections:

\begin{enumerate}
  \item slightly downstream of the inlet,
  \item at the midpoint of the straight inlet,
  \item approximately one channel width upstream of the junction.
\end{enumerate}

For each section, the axial velocity \(u_\parallel\) is compared with

\[
  u_{\mathrm{Pois}}(s)
  =
  \frac{3}{2}U_{\mathrm{mean}}
  \left[
  1-\left(\frac{2s}{W}\right)^2
  \right].
\]

For the production 50/50 case, the flux errors were approximately
\(-2.778\times10^{-4}\), \(-2.771\times10^{-4}\), and
\(-2.818\times10^{-4}\). The maximum transverse-to-mean inlet velocity ratio
was about \(9.98\times10^{-4}\). These values show that the straight inlet
remains close to the imposed Poiseuille profile and that the previous flux
collapse failure has been eliminated.

\subsection{Version 1 Prescribed-Split Library}

The full Version 1 library consists of 19 prescribed split cases:

\[
  f_L \in
  \{0.05, 0.10, 0.15,\ldots,0.95\}.
\]

The endpoints \(0.00\) and \(1.00\) are intentionally excluded because
zero-flow outlet treatment may require a separate boundary-condition analysis.

Each case is stored under:

\begin{verbatim}
outputs/physics/junction_cfd/solutions/split_0pXX/
\end{verbatim}

Each case contains solution fields, metadata, flux report, pressure field,
velocity magnitude, velocity vectors, dense streamlines, inlet-seeded
streamlines, and inlet Poiseuille diagnostics.

\begin{center}
\small
\begin{longtable}{@{}lrrrrr@{}}
\toprule
Case & Requested \(f_L\) & Measured \(f_L\) & Residual & Pressure range & Seeds L/R/O \\
\midrule
\endhead
\texttt{split\_0p05} & 0.05 & 0.0500 & \(9.096{\times}10^{-17}\) & 37.00166 & 0/31/0 \\
\texttt{split\_0p10} & 0.10 & 0.1000 & \(7.893{\times}10^{-17}\) & 36.00268 & 0/31/0 \\
\texttt{split\_0p15} & 0.15 & 0.1500 & \(6.690{\times}10^{-17}\) & 35.00370 & 0/31/0 \\
\texttt{split\_0p20} & 0.20 & 0.2000 & \(5.487{\times}10^{-17}\) & 34.00472 & 1/28/2 \\
\texttt{split\_0p25} & 0.25 & 0.2500 & \(4.284{\times}10^{-17}\) & 33.00574 & 4/26/1 \\
\texttt{split\_0p30} & 0.30 & 0.3000 & \(3.081{\times}10^{-17}\) & 32.00676 & 6/23/2 \\
\texttt{split\_0p35} & 0.35 & 0.3500 & \(1.878{\times}10^{-17}\) & 31.00778 & 8/21/2 \\
\texttt{split\_0p40} & 0.40 & 0.4000 & \(6.754{\times}10^{-18}\) & 30.00879 & 11/19/1 \\
\texttt{split\_0p45} & 0.45 & 0.4500 & \(-5.276{\times}10^{-18}\) & 29.00981 & 13/17/1 \\
\texttt{split\_0p50} & 0.50 & 0.5000 & \(-1.730{\times}10^{-17}\) & 28.70440 & 15/15/1 \\
\texttt{split\_0p55} & 0.55 & 0.5500 & \(-2.933{\times}10^{-17}\) & 30.05921 & 17/13/1 \\
\texttt{split\_0p60} & 0.60 & 0.6000 & \(-4.136{\times}10^{-17}\) & 31.41402 & 19/10/2 \\
\texttt{split\_0p65} & 0.65 & 0.6500 & \(-5.339{\times}10^{-17}\) & 32.76883 & 21/8/2 \\
\texttt{split\_0p70} & 0.70 & 0.7000 & \(-6.542{\times}10^{-17}\) & 34.12364 & 23/6/2 \\
\texttt{split\_0p75} & 0.75 & 0.7500 & \(-7.745{\times}10^{-17}\) & 35.47845 & 26/4/1 \\
\texttt{split\_0p80} & 0.80 & 0.8000 & \(-8.948{\times}10^{-17}\) & 36.83327 & 28/1/2 \\
\texttt{split\_0p85} & 0.85 & 0.8500 & \(-1.015{\times}10^{-16}\) & 38.18808 & 30/0/1 \\
\texttt{split\_0p90} & 0.90 & 0.9000 & \(-1.135{\times}10^{-16}\) & 39.54289 & 31/0/0 \\
\texttt{split\_0p95} & 0.95 & 0.9500 & \(-1.256{\times}10^{-16}\) & 40.89770 & 31/0/0 \\
\bottomrule
\end{longtable}
\end{center}

All 19 cases passed validation. The maximum absolute mass residual was
\(1.256\times10^{-16}~\mathrm{m^2/s}\). The maximum requested-versus-measured
left-fraction error was \(1.735\times10^{-11}\). Cross-case checks confirmed
that left outlet flux increases monotonically with \(f_L\), right outlet flux
decreases monotonically, and the measured split tracks the requested split.

\begin{figure}[H]
  \centering
  \includegraphics[width=0.95\linewidth,height=0.80\textheight,keepaspectratio]{../outputs/physics/junction_cfd/solutions/library_figures/library_summary_metrics.png}
  \caption{Summary metrics across the 19 prescribed split cases. Fluxes provide the quantitative split validation; streamline counts are qualitative diagnostics.}
  \label{fig:cfd-library-summary}
\end{figure}

\begin{figure}[H]
  \centering
  \includegraphics[width=0.95\linewidth,height=0.80\textheight,keepaspectratio]{../outputs/physics/junction_cfd/solutions/library_figures/library_representative_montage.png}
  \caption{Representative velocity magnitude and inlet-seeded streamlines for selected prescribed split cases \(f_L=0.05,0.25,0.50,0.75,0.95\).}
  \label{fig:cfd-library-montage}
\end{figure}

\subsection{Execution}

Build the production mesh:

\begin{verbatim}
.\.venv\Scripts\python.exe -m scripts.physics.build_junction_mesh --overwrite
\end{verbatim}

Run the default validated 50/50 case:

\begin{verbatim}
.\.venv\Scripts\python.exe -m scripts.physics.solve_junction_cfd --overwrite
\end{verbatim}

Run one explicit prescribed split:

\begin{verbatim}
.\.venv\Scripts\python.exe -m scripts.physics.solve_junction_cfd ^
  --left-fraction 0.30 --overwrite
\end{verbatim}

Generate the full Version 1 library:

\begin{verbatim}
.\.venv\Scripts\python.exe -m scripts.physics.generate_velocity_library --overwrite
\end{verbatim}

\subsection{Frozen Version 1 Artifacts}

The frozen baseline lives under:

\begin{verbatim}
outputs/physics/junction_cfd/version_1/
\end{verbatim}

The concise Version 1 manifest is:

\begin{verbatim}
outputs/physics/junction_cfd/version_1/CFD_VERSION_1.md
\end{verbatim}

These artifacts record the validated scientific formulation and should not be
overwritten by future library generation or transformer-facing experiments.
